% Generated draft: requires mathematical and source review.
\documentclass[11pt]{article}
\usepackage[margin=25mm]{geometry}
\usepackage{fontspec}
\setmainfont{Latin Modern Roman}
\usepackage{amsmath,amssymb,mathtools,unicode-math}
\setmathfont{Latin Modern Math}
\usepackage{enumitem,needspace,xurl,hyperref}
\setlength{\parindent}{0pt}
\setlength{\parskip}{5pt}
\setlength{\emergencystretch}{3em}
\newcommand{\N}{\mathbb{N}}
\newcommand{\Z}{\mathbb{Z}}
\newcommand{\Q}{\mathbb{Q}}
\newcommand{\R}{\mathbb{R}}
\newcommand{\C}{\mathbb{C}}
\newcommand{\F}{\mathbb{F}}
\newcommand{\A}{\mathbb{A}}
\newcommand{\PP}{\mathbb{P}}
\newcommand{\E}{\mathbb{E}}
\newcommand{\eps}{\varepsilon}
\newcommand{\dd}{\,\mathrm{d}}
\begin{document}
\section*{UnsolvedMath ID: 30004364}
\textbf{Source problem code:} catalogue entry
\subsection{Definitions and mathematical statement}
The supplied record does not specify the base field, the permitted quivers, or the precise convention for \(\operatorname{rep}Q\). The following expresses its recovery criterion conditionally, in a setting where representations have finite-dimensional vertex spaces and the generic data described below exist; these qualifications are not additional hypotheses supplied by the source.

Write \(Q_0\) and \(Q_1\) for the vertices and arrows of a quiver \(Q\). Over a fixed field \(k\), a representation \(M\) assigns a vector space \(M_i\) to each vertex \(i\) and a linear map \(M_a:M_i\to M_j\) to each arrow \(a:i\to j\). An endomorphism \(f=(f_i)_{i\in Q_0}\) satisfies
\[
 f_jM_a=M_af_i\qquad(a:i\to j).
\]
It is nilpotent if \(f^r=0\) for some integer \(r\geq1\). For finite-dimensional \(M_i\), let \(\lambda_i(f)\) be the partition of \(\dim_k M_i\) consisting of the sizes of the nilpotent Jordan blocks of \(f_i\); the zero space has the empty partition. Put
\[
 \operatorname{NilEnd}(M)=\{f\in\operatorname{End}(M):f^r=0\text{ for some }r\geq1\}.
\]
Whenever the intended algebraic parameter space \(\operatorname{NilEnd}(M)\) has a dense open subset on which the tuple \((\lambda_i(f))_{i\in Q_0}\) is constant, denote this generic tuple by
\[
 \operatorname{GenJF}(M)=(\lambda_i(M))_{i\in Q_0}.
\]
The source does not specify the geometric setting or hypotheses guaranteeing this generic value.

An additive subcategory has a zero object and finite direct sums with compatible additive morphism structure. Whether the subcategories here must additionally be full, closed under isomorphism, or closed under direct summands is not specified.

\textbf{Problem.} Characterize the additive subcategories \(\mathcal C\subseteq\operatorname{rep}Q\) for which generic nilpotent Jordan data are defined for every object and determine its isomorphism class within \(\mathcal C\). Explicitly, the recovery criterion suggested by the supplied wording is
\[
 \forall M,N\in\operatorname{Ob}(\mathcal C),\qquad
 \operatorname{GenJF}(M)=\operatorname{GenJF}(N)
 \ \Longrightarrow\ M\cong N.
\]
Equality here means equality of the vertex-indexed partitions. This criterion interprets recovery as uniqueness among objects of \(\mathcal C\); the source does not explicitly distinguish this from uniqueness in the ambient category or from a stronger canonical reconstruction requirement. A definitive formulation requires resolving these conventions and the missing ambient hypotheses.
\subsection{Short English statement}
Classify the additive subcategories of a quiver representation category in which the vertexwise Jordan block sizes of a generic nilpotent endomorphism determine each representation up to isomorphism. The field, allowed quivers, subcategory conventions, and precise scope of recovery remain unspecified in the supplied record.
\subsection{Sources}
\begin{itemize}
\item Supplied local UnsolvedMath registry record, permanent ID \(30004364\), title \emph{Jordan-Recoverable Quiver Subcategories}, exact source problem code \texttt{catalogue entry}, in the Oberwolfach Reports Open Problems collection. The record lists proposed year \(2020\) and supplies no individual proposer.
\item Supplied report citation: \emph{Representation Theory of Quivers and Finite Dimensional Algebras} (2021), Oberwolfach Reports, pp.\ 143--230, DOI \texttt{10.4171/owr/2020/3}. \url{https://doi.org/10.4171/owr/2020/3}. The supplied background describes these reports as documenting workshops and their open-problem discussions. No problem-specific page or speaker attribution is supplied.
\item Supplied literature reference: Dequ\^ene, \emph{Canonically Jordan recoverable categories}. \url{https://arxiv.org/abs/2308.16626}.
\item Literature assessment dated exactly \texttt{2026-08-22}: the supplied curation labels the problem open and reports partial constructions and classifications for special quivers and subcategories, with no general characterization of all Jordan recoverable additive subcategories. This is supplied dated curation, not an independently verified assessment; its time-sensitive status and cited claims require independent verification.
\item The registry labels its statement \texttt{corrected\_verified} and reports manual semantic, typographic, and LaTeX review. These are supplied editorial metadata, not independent verification of the mathematical formulation or literature status in this draft.
\end{itemize}
\end{document}
